\documentclass[11pt,a4paper]{article}

%% ---------------------------------------------------------------------------
%% On the Extremal Number of Maximal Convex-Position Subsets
%% of a Planar Point Set
%%
%% Compile with:  pdflatex main.tex   (twice, for the bibliography)
%% The numbers in this paper are produced by
%%     experiments/reproduce_paper_numbers.py
%% which writes results/paper_numbers.json.
%% ---------------------------------------------------------------------------

\usepackage[margin=2.6cm]{geometry}
\usepackage[T1]{fontenc}
\usepackage[utf8]{inputenc}
\usepackage{amsmath,amssymb,amsthm}
\usepackage{booktabs}
\usepackage{enumitem}
\usepackage[hidelinks]{hyperref}

\hypersetup{
  pdftitle={On the Extremal Number of Maximal Convex-Position Subsets of a Planar Point Set},
  pdfauthor={Rajveersinh Pardeshi},
}

\newtheorem{theorem}{Theorem}[section]
\newtheorem{lemma}[theorem]{Lemma}
\newtheorem{proposition}[theorem]{Proposition}
\newtheorem{corollary}[theorem]{Corollary}
\theoremstyle{definition}
\newtheorem{definition}[theorem]{Definition}
\newtheorem{remark}[theorem]{Remark}
\newtheorem{conjecture}[theorem]{Conjecture}
\newtheorem{openproblem}[theorem]{Open problem}

\newcommand{\conv}{\operatorname{conv}}
\newcommand{\vtx}[1]{\operatorname{vert}(#1)}
\newcommand{\M}{\mathcal{M}}
\newcommand{\CP}{\mathcal{C}}
\newcommand{\HP}{\mathcal{H}}
\newcommand{\intf}{\mathrm{int}}
\newcommand{\hull}{\mathrm{Hull}}

\title{\bf On the Extremal Number of Maximal\\
       Convex-Position Subsets of a Planar Point Set}
\author{Rajveersinh Pardeshi}
\date{October 2026}

\begin{document}
\maketitle

\begin{abstract}
A subset $S$ of a finite point set $P\subset\mathbb{R}^2$ in general position is in
\emph{convex position} if every point of $S$ is a vertex of $\conv(S)$, and is a
\emph{maximal convex subset} if no proper superset of $S$ inside $P$ is in convex
position.  Writing $\M(P)$ for the family of maximal convex subsets of $P$ and
$f(n)=\max_{|P|=n}|\M(P)|$, we study the growth of $f$.

We compute $f$ exactly for $n\le 9$ by an exhaustive scan of a complete database
of order types of planar point sets:
\[
  f(3)=1,\quad f(4)=4,\quad f(5)=7,\quad f(6)=11,\quad
  f(7)=20,\quad f(8)=38,\quad f(9)=62 .
\]
Our main structural result is a duality theorem: the map
$\psi(S)=P\setminus\conv(S)$ is injective on the \emph{entire} family of
convex-position subsets --- maximality is not needed --- and it identifies that
family with the lattice of \emph{hull-closed} subsets $W\subseteq P$ satisfying
$W=P\cap\conv(W)$.  In particular two distinct maximal convex subsets never have
the same hull as a subset of $P$.  We also prove a kill lemma: if $S\in\M(P)$ and
$p\in P\setminus S$ lies outside $\conv(S)$, then $p$ destroys at least one
vertex of $S$.  We further show that $|\M(P)|$ is extraordinarily sensitive to
the order type --- on fixed $14$ points it ranges over $34,\dots,121$ as a single
rational parameter varies --- and we report, with corrections, two of our own
initial guesses that the data refuted.  We conjecture $f(n)=2^{\Theta(n)}$ and
state precisely what remains open.
\end{abstract}

\section{Introduction}
\label{sec:intro}

\subsection{The problem}

Let $P\subset\mathbb{R}^2$ be a finite set of $n$ points in \emph{general
position}: no three points are collinear.  A subset $S\subseteq P$ with $|S|\ge 3$
is \emph{in convex position} if each of its points is a vertex of $\conv(S)$,
equivalently if $S$ is the vertex set of a convex polygon.  A convex-position
subset is \emph{maximal} if no $S'\subsetneq P$ properly containing $S$ is in
convex position.  Put
\[
  \M(P)=\bigl\{S\subseteq P:\ S \text{ is a maximal convex-position subset of } P\bigr\}
\]
and
\[
  f(n)=\max\bigl\{|\M(P)|:\ P\subset\mathbb{R}^2,\ |P|=n,\ \text{\(P\) in general position}\bigr\}.
\]

The question \emph{how many convex polygons does a point set determine} has a long
history~\cite{rote91,rote92,mitchell95,huemer2022}: those works count \emph{all}
convex subsets, or count them for a \emph{given} configuration.  Two neighbouring
extremal questions are different again.  The \emph{minimum} number of convex
subsets of $n$ points is Erd{\\H{o}}s Problem~838~\cite{erdos838}, on the order
$n^{\Theta(\log n)}$; exact values for $n\le 11$ were recently computed
by Savova~\cite{savova2026}.  The number of \emph{empty} convex polygons (holes)
is controlled by~\cite{baranyvaltr04,pinchasi2006,bae2022,suk2017}.

Our question is a third one.  We \emph{maximise} the number of
\emph{inclusion-maximal} convex subsets over all configurations.  Maximising the
number of convex $k$-gons for a fixed $k$ is trivial --- it is $\binom{n}{k}$,
attained when $P$ is in convex position --- so the interesting object is the
number of convex subsets that cannot be enlarged.  This is a purely order-type
invariant (Section~\ref{sec:prelim}) encoding the local structure of $P$, and it
records information that neither the total count of convex subsets nor the count
of convex $k$-gons captures.

We are not aware of any prior work on $f$.  See Section~\ref{sec:related} for a
discussion of the literature we searched and of the terminology we checked.

\subsection{Why the quantity is non-trivial}

By Lemma~\ref{lem:downward} the family $\M(P)$ is an \emph{antichain} in $2^P$: if
$S\subsetneq T\subseteq P$ and $T$ is in convex position then so is $S$, so $S$
is not maximal.  Hence Sperner's theorem applies and
\begin{equation}
  |\M(P)|\ \le\ \binom{n}{\lfloor n/2\rfloor}\ =\ 2^{\Theta(n)} .
  \label{eq:sperner}
\end{equation}
The bound \eqref{eq:sperner} is weak, and we do not improve it in general.  What
we obtain is (i) the exact values of $f$ for $n\le 9$; (ii) a duality theorem that
governs the structure of the whole family of convex-position subsets; (iii) a
local geometric lemma; and (iv) explicit configurations showing that $f$ grows
quickly and is extremely sensitive to the order type.

\subsection{Contributions}

\begin{enumerate}[leftmargin=*]
\item \textbf{Exact values (Theorem~\ref{thm:f}).} $f(n)$ for $n=3,\dots,9$,
  computed exhaustively over \emph{all} order types of $n$ points.
\item \textbf{Duality (Theorem~\ref{thm:duality}).} $\psi(S)=P\setminus\conv(S)$
  is injective on the entire family of convex-position subsets of size $\ge 3$,
  and identifies it with the family of hull-closed subsets.  Maximality is not
  used in the proof.
\item \textbf{Kill lemma (Lemma~\ref{lem:kill}).} If $S\in\M(P)$ and
  $p\in P\setminus S$ lies outside $\conv(S)$, then $p$ destroys at least one
  vertex of $S$.
\item \textbf{Layer profile (Proposition~\ref{prop:layer}).} The maximum of
  $|\M(P)|$ over configurations with $h$ hull vertices is attained at $h=3$ or
  $4$, and collapses to $1$ when $P$ is in convex position.
\item \textbf{Sensitivity (Section~\ref{sec:twin}).} On fixed $14$ points,
  $|\M(P)|$ ranges over $34,\dots,121$ as a single rational separation parameter
  varies.  We also record that our explicit constructions are \emph{weaker} than
  exhaustive search at small $n$, so they support no asymptotic lower bound.
\end{enumerate}

\subsection{Organisation}

Section~\ref{sec:prelim} fixes notation.  Section~\ref{sec:duality} proves the
duality theorem.  Section~\ref{sec:local} gives the local structural results.
Section~\ref{sec:related} discusses related work and novelty.  Section~\ref{sec:comp}
describes the computational method in enough detail to reproduce it, including two
bugs we found and fixed, and gives complexity bounds.  Section~\ref{sec:values}
reports the exact values.  Section~\ref{sec:twin} gives constructions and the
falsification experiments, including two of our own conjectures that our data
refuted.  Section~\ref{sec:open} states what is open.

\section{Preliminaries}
\label{sec:prelim}

Throughout, $P$ is a finite set of $n$ points in general position in
$\mathbb{R}^2$.  For $S\subseteq P$ let $\conv(S)$ be the closed convex hull and
$\vert(S)$ the set of vertices of $\conv(S)$.

\begin{definition}[convex position]
A set $S\subseteq P$ with $|S|\ge 3$ is \emph{in convex position} if
$S=\vert(S)$, that is, every point of $S$ is an extreme point of $\conv(S)$.
Equivalently, $S$ is the vertex set of a convex $|S|$-gon.
\end{definition}

\begin{definition}[maximal convex]
A convex-position subset $S\subseteq P$ is \emph{maximal convex} if for every
$p\in P\setminus S$ the set $S\cup\{p\}$ is not in convex position.
\end{definition}

\begin{definition}
$\CP(P)$ is the family of convex-position subsets of $P$ of size at least $3$;
$\M(P)$ is the family of maximal members of $\CP(P)$; and $f(n)$ is the maximum
of $|\M(P)|$ over all $n$-point general-position sets.
\end{definition}

\begin{definition}[forbidden quadruple]
A $4$-subset $\{a,b,c,d\}\subset P$ is \emph{forbidden} if one of its points lies
strictly inside the triangle spanned by the other three.  In general position at
most one point can be interior, so the witness is unique.
\end{definition}

\begin{definition}[hull-closed]
A subset $W\subseteq P$ with $|W|\ge 3$ is \emph{hull-closed} if
$W=P\cap\conv(W)$: every point of $P$ lying in $\conv(W)$ already belongs to $W$.
\end{definition}

\begin{definition}[order type]
The \emph{order type} of $P$ is the isomorphism class of the map
$\{i,j,k\}\mapsto\operatorname{sgn}\operatorname{orient}(p_i,p_j,p_k)$.  Every
quantity we study depends only on the order type of $P$.
\end{definition}

\begin{lemma}[downward closure]
\label{lem:downward}
If $S\in\CP(P)$ and $T\subseteq S$ with $|T|\ge 3$ then $T\in\CP(P)$.
Consequently $\CP(P)$ is a simplicial complex on vertex set $P$ and $\M(P)$ is the
set of its facets; in particular $\M(P)$ is an antichain.
\end{lemma}

\begin{proof}
If $T\subseteq S$ and some $t\in T$ were in $\int(\conv(T))$, then $t$ would lie
in $\int(\conv(S))$ as well, contradicting $S\in\CP(P)$.  A subset of a set of
points in convex position is again in convex position, so $\CP(P)$ is a simplicial
complex and $\M(P)$ is its set of facets.  Facets of a simplicial complex are
pairwise incomparable, so $\M(P)$ is an antichain.
\end{proof}

\begin{proposition}[Sperner bound]
\label{prop:sperner}
$|\M(P)|\le\binom{n}{\lfloor n/2\rfloor}$.
\end{proposition}

\begin{proof}
By Lemma~\ref{lem:downward}, $\M(P)$ is an antichain in $2^P$.  Apply Sperner's
theorem.
\end{proof}

\section{The duality theorem}
\label{sec:duality}

\begin{theorem}[duality]
\label{thm:duality}
Let $P$ be in general position and let $S,T\subseteq P$ be in convex position
with $|S|,|T|\ge 3$.  Then
\[
  P\setminus\conv(S)=P\setminus\conv(T)\quad\Longrightarrow\quad S=T.
\]
\end{theorem}

\begin{proof}
Put $W=P\cap\conv(S)=P\cap\conv(T)$.

Since $S\subseteq\conv(S)$ we get $S\subseteq P\cap\conv(S)=W$, and likewise
$T\subseteq W$.  Because $S$ is in convex position, $\vert(S)=S$, so
$S\subseteq W\subseteq\conv(T)$.  Using $S=\vert(S)$ once more,
\[
  \conv(S)=\conv\bigl(\vert(S)\bigr)\ \subseteq\ \conv(T).
\]
By symmetry $\conv(T)\subseteq\conv(S)$.  Therefore $\conv(S)=\conv(T)$ and
\[
  S=\vert(S)=\vert(T)=T.
\]
\end{proof}

\begin{remark}
Maximality plays no role in the proof of Theorem~\ref{thm:duality}: the statement
holds for $\CP(P)$ itself, not merely for $\M(P)$.
\end{remark}

\begin{corollary}[identification with hull-closed subsets]
\label{cor:bijection}
The map $\psi:S\mapsto P\setminus\conv(S)$ is a bijection from $\CP(P)$ onto
$\HP(P)=\{W\subseteq P:\ |W|\ge 3,\ W=P\cap\conv(W)\}$.  In particular
$|\CP(P)|=|\HP(P)|$.
\end{corollary}

\begin{proof}
Injectivity is Theorem~\ref{thm:duality}.  For surjectivity, let $W$ be
hull-closed and set $S=\vert(W)$.  Then $S\subseteq W\subseteq\conv(S)$, so
$\conv(S)=\conv(W)$ and
$P\cap\conv(S)=P\cap\conv(W)=W$, i.e.\ $\psi(S)=W$.  Also $S$ is in convex
position by construction, and $S\subseteq W=P\cap\conv(W)$ with
$|W|\ge 3$ forces $|S|\ge 3$ (otherwise $W$ would be a segment or a point and,
by general position, $W\cap P=W$ would have $|W|\le 2$).
\end{proof}

\begin{corollary}[distinct hulls]
\label{cor:distinct}
If $S,T\in\M(P)$ and $S\ne T$, then $P\cap\conv(S)\ne P\cap\conv(T)$.  In
particular $P\setminus\conv(S)\ne P\setminus\conv(T)$: distinct maximal convex
subsets are always distinguished by their hull.
\end{corollary}

\section{Local structure}
\label{sec:local}

\subsection{The kill lemma}

\begin{lemma}[kill lemma]
\label{lem:kill}
Let $S\in\M(P)$, let $p\in P\setminus S$, and suppose $p\notin\conv(S)$.  Then
there is $v\in S$ that lies strictly in $\int\bigl(\conv((S\setminus\{v\})\cup\{p\})\bigr)$.
\end{lemma}

\begin{proof}
Since $S$ is in convex position, $\vert(S)=S$.  Every point of $S$ is therefore a
vertex of $\conv(S)$.

Suppose, for a contradiction, that the conclusion fails.  Then for every $v\in S$
the point $v$ remains a vertex of $\conv((S\setminus\{v\})\cup\{p\})$; in
particular $v$ is still a vertex of $\conv(S\cup\{p\})$.

Fix $v\in S$.  Because $S\cup\{p\}$ consists of $S$ together with the single
extra point $p$, and $v$ survives, $v$ is a vertex of $\conv(S\cup\{p\})$.  The
same holds for every $v\in S$, and $p$ is trivially a vertex of
$\conv(S\cup\{p\})$.  Hence
\[
  \vert(S\cup\{p\})=S\cup\{p\},\qquad |S\cup\{p\}|=|S|+1,
\]
so $S\cup\{p\}$ is in convex position, contradicting $S\in\M(P)$.
\end{proof}

\begin{remark}
Lemma~\ref{lem:kill} says that adding a point outside the hull of a maximal
convex subset always costs at least one vertex.  Empirically it usually costs
exactly one: over all pairs $(S,p)$ with $n\le 7$ and $p\notin\conv(S)$, the
numbers of destroyed vertices $1,2,3,4$ occur with frequencies
$3438,828,46,1$, with no violations of the lemma among $6126$ pairs.
\end{remark}

\subsection{Hull layers}

\begin{lemma}
\label{lem:convexpos}
If $P$ is in convex position then $\M(P)=\{P\}$ and $|\M(P)|=1$.
\end{lemma}

\begin{proof}
Every subset of a convex-position set of size at least $3$ is in convex position
(Lemma~\ref{lem:downward}), so $P$ itself is maximal and nothing smaller is.
\end{proof}

\begin{proposition}[layer profile]
\label{prop:layer}
Let $g(n,h)$ be the maximum of $|\M(Q)|$ over general-position $n$-point sets
$Q$ with $|\hull(Q)|=h$ hull vertices.  For $n\le 8$ the maximum of $g(n,h)$ over
$h$ is attained at $h\in\{3,4\}$, and $g(n,n)=1$.
\end{proposition}

\begin{proof}
The statement for $n\le 8$ is Table~\ref{tab:layer}, obtained by an exhaustive scan
over all order types; $g(n,n)=1$ is Lemma~\ref{lem:convexpos}.
\end{proof}

\section{Related work and novelty}
\label{sec:related}

\subsection{What is known}

Counting convex subsets of planar point sets is well studied.
Rote, Woeginger, Zhu and Wang~\cite{rote91} count convex $k$-gons; Rote and
Woeginger~\cite{rote92} and Mitchell, Rote, Sundaram and
Woeginger~\cite{mitchell95} give algorithms computing, in $O(n^3)$ time, the total
number of convex polygons determined by a given $n$-point set, and Mitchell et
al.\ also compute, for a given point $u$, the number of convex $k$-gons containing
$u$.  Huemer, Oliveros, P\'erez-Lantero, Torra and Vogtenhuber~\cite{huemer2022}
prove linear identities among the counts $X_{k,\ell}(S)$ of convex $k$-gons with
$\ell$ interior points.

For the \emph{minimum} number of convex subsets, Erd{\\H{o}}s Problem~838~\cite{erdos838}
asks for the order of magnitude of $\min_{|P|=n}|\CP(P)|$ (with singletons and pairs
counted), and Savova~\cite{savova2026} reports exact values for $n\le 11$ together
with the bound $\limsup\log_2 f(N)/(\log_2 N)^2\le \tfrac12$.

\paragraph{What is different here.}

Three distinctions matter.

\begin{enumerate}[leftmargin=*]
\item \emph{Maximum versus minimum.} Erd{\\H{o}}s Problem~838 and~\cite{savova2026}
  concern $\min_P|\CP(P)|$.  We concern $\max_P|\M(P)|$, where $\M(P)$ counts only
  the inclusion-maximal convex subsets.
\item \emph{Maximal versus all.} The papers~\cite{rote91,rote92,mitchell95} count
  \emph{all} convex subsets of a \emph{fixed} $P$.  Counting only the maximal ones
  is a genuinely different functional: for $P$ in convex position,
  $|\CP(P)|=2^n-1-n-\binom n2$ while $|\M(P)|=1$.
\item \emph{Order types.} Since $|\M(P)|$ is order-type invariant, the whole problem
  reduces to a finite computation for each $n$.  We use the
  Aichholzer--Aurenhammer--Krasser database~\cite{aak2002,aak2001,aak2007}, which
  stores one realisation per order type for $n\le 10$.
\end{enumerate}

\subsection{Terminology checked}

We searched for prior work under many phrasings, including
``maximal convex position subset'', ``inclusion-maximal convex subset'',
``convex position complex'', ``facets of the convex position complex'',
``maximal convex polygon determined by a point set'', ``saturated convex
position'', and ``$\varepsilon$-nets in convex position''.

The quantity we study is the number of facets of the \emph{convex position
complex}: the simplicial complex whose faces are the convex-position subsets of
$P$.  Work on related complexes for convex \emph{bodies} --- notably on
$\varepsilon$-nets for convex bodies in convex position~\cite{hmp2020} --- and on
generalised convexity notions such as stair-convexity~\cite{bukhmnk11} does arise
from a complex of convex-position subsets.  Those works concern the topology and
the optimisation aspects in terms of $d$, $n$ and $\varepsilon$; they do not ask
for the extremal number of facets of the convex position complex of a planar point
set as a function of $n$, which is our question.  We did not locate a paper that
treats that extremal problem.

We also note that the phrase ``maximal empty convex polygon'' occurs in the holes
literature~\cite{baranyvaltr04,pinchasi2006}; maximality there refers to maximal
families of disjoint holes, a different notion from our inclusion-maximality among
convex-position subsets.

\subsection{Honest novelty statement}

We did \emph{not} find a prior result determining $f(n)$ for any $n$, or bounding
it non-trivially.  We therefore do not claim to be the first to consider the
extremal count of maximal convex-position subsets; we claim that we did not find
such a result, and that the values, the duality theorem and the constructions
reported here are new to us.  A reader with access to MathSciNet should re-check
this claim; our search was conducted over zbMATH Open, Crossref, OpenAlex, arXiv
listing and abstract search, and MathOverflow.

\section{Computational method}
\label{sec:comp}

This section is deliberately explicit so that the numbers in
Section~\ref{sec:values} can be checked independently.

\subsection{Exactness}

Every geometric predicate is decided by an orientation sign of \emph{rational}
coordinates (Python ``fractions.Fraction``); no floating-point value is ever used
to decide an orientation.  A configuration is used only through its chirotope, so
all arithmetic is exact.

\subsection{Counting $\M(P)$ two ways}

We implement two independent routines and cross-check them.

\paragraph{Brute force.} Enumerate all subsets of size $\ge 3$, test convex
position by computing the hull with Andrew's monotone chain, then delete those
properly contained in another convex-position subset.  Cost $O(2^n n\log n)$,
usable to $n\approx 12$.

\paragraph{Forbidden-quadruple enumeration.} By the characterisation
\[
  S\in\CP(P)\quad\Longleftrightarrow\quad S\ \text{contains no forbidden quadruple},
\]
$\CP(P)$ is the independent-set family of the $4$-uniform hypergraph $H_P$ of
forbidden quadruples, and $\M(P)$ is the family of \emph{maximal} independent sets
of $H_P$.  We enumerate maximal independent sets output-sensitively: at a node
with chosen set $S$ and candidate set $R$ (exactly the vertices still compatible
with $S$), if $R=\varnothing$ then $S$ is maximal and we record it; otherwise we
branch on each $v\in R$.  Every leaf is a maximal independent set and every
maximal independent set labels exactly one leaf, so the cost is exponential in
$|\M(P)|$ rather than in $2^n$.  A node budget makes truncation raise an error
rather than silently return a truncated answer.  Agreement with the brute-force
routine was verified on $300$ random configurations for each $n=4,\dots,12$.

\subsection{Order-type database}

Because $|\M(P)|$ depends only on the order type, one realisation per order type
suffices.  We use the Aichholzer--Aurenhammer--Krasser
database~\cite{aak2002,aak2001,aak2007}.  Its documented format (section B of the
accompanying ``readme.txt'') is: no header; for each order type the coordinates are
stored in the order $x_1,y_1,x_2,y_2,\dots,x_n,y_n$; coordinates are
\emph{unsigned} integers, one byte for $n\le 8$ and two bytes (low byte first) for
$n=9,10$.  Our reader reproduces the classical counts
$1,2,3,16,135,3315,158817$ for $n=3,\dots,9$, and we verify that every stored
realisation is in general position with distinct coordinates
(\texttt{tests/test\_aak\_database.py}).

\subsection{An independent enumerator, and two bugs it caught}
\label{sec:bugs}

To guard against trusting a single tool, we also implemented a complete
enumeration of order types by \emph{incremental insertion}: every order type on $n$
points restricts to one on $n-1$ points, and each extension is determined by the
cell of the arrangement of pair-lines occupied by the new point.  Cells are
enumerated by a vertical sweep --- the critical abscissae (where two non-parallel
lines meet) split the plane into slabs, and inside a slab each gap between
consecutive lines contains exactly one cell --- using rational sample points and a
generic rational rotation to remove vertical lines.

This independent route found two bugs, which we report openly.

\begin{enumerate}[leftmargin=*]
\item Our sweep initially used $x=(b_2c_1-b_1c_2)/\det$ instead of
  $x=(b_1c_2-b_2c_1)/\det$ for the intersection abscissa of two lines
  $a_ix+b_iy+c_i=0$, so the wrong slabs were sampled.  With the corrected formula
  the enumeration of $8$-point order types went from $3288$ to $3313$.
\item Even then two types out of $3315$ were missing, because the sampler excluded
  vertical pair-lines from the gap scan.  Applying a generic rational rotation
  before the sweep, and mapping sample points back afterwards, preserves the cell
  combinatorics; this is implemented but we did not re-run the full $n=8$
  enumeration afterwards.
\end{enumerate}

We therefore do \emph{not} claim our own enumerator is complete at $n=8$: it
reaches $3313$ of $3315$.  All exact values in Section~\ref{sec:values} come from
the database scan, which is complete by construction.

\section{Exact values of $f$}
\label{sec:values}

\begin{theorem}
\label{thm:f}
For $n=3,\dots,9$,
\[
  f(3)=1,\quad f(4)=4,\quad f(5)=7,\quad f(6)=11,\quad
  f(7)=20,\quad f(8)=38,\quad f(9)=62.
\]
\end{theorem}

\begin{proof}
Computed by exhaustive scan of all order types of $n$ points (Table~\ref{tab:f}).
The database contains exactly one realisation per order type and $|\M(P)|$ is
order-type invariant, so the scan is exhaustive; since the set of order types is
finite and fully known, the maximum found is $f(n)$.
\end{proof}

\begin{table}[ht]
\centering
\caption{$f(n)$ and the number of order types scanned.}
\label{tab:f}
\begin{tabular}{r r r l}
\toprule
$n$ & order types & $f(n)$ & sizes of the members of a maximising $\M(P)$ \\
\midrule
3 & 1       & 1  & one set of size $3$ \\
4 & 2       & 4  & four sets of size $3$ \\
5 & 3       & 7  & six of size $3$, one of size $4$ \\
6 & 16      & 11 & two of size $3$, nine of size $4$ \\
7 & 135     & 20 & one of size $3$, nineteen of size $4$ \\
8 & 3315    & 38 & thirty-eight of size $4$ \\
9 & 158817  & 62 & see \texttt{results/f\_aak.json} \\
\bottomrule
\end{tabular}
\end{table}

\begin{table}[ht]
\centering
\caption{$\max|\M(P)|$ as a function of $n$ and the hull size $h$, for $n\le 8$.
The optimum sits at $h=3$ or $h=4$; configurations in convex position ($h=n$) give
$|\M(P)|=1$ by Lemma~\ref{lem:convexpos}.}
\label{tab:layer}
\begin{tabular}{c r r r r r r}
\toprule
$n \backslash h$ & 3 & 4 & 5 & 6 & 7 & 8 \\
\midrule
4 & 4 & 1 &   &   &   &   \\
5 & 7 & 4 & 1 &   &   &   \\
6 & 11 & 11 & 6 & 1 &  &   \\
7 & 20 & 20 & 12 & 6 & 1 &  \\
8 & 35 & \textbf{38} & 26 & 16 & 8 & 1 \\
\bottomrule
\end{tabular}
\end{table}

Two structural observations, both exact for $n\le 8$:

\begin{itemize}[leftmargin=*]
\item At the optimum for $n=6$ and $n=8$, \emph{every} maximal convex subset has
  the same size (namely $n/2$) and \emph{every} point of $P$ lies in the same
  number of maximal subsets.  The extremal configurations are thus perfectly
  balanced.
\item In the $n=8$ maximiser the $38$ maximal subsets are all $4$-subsets, so
  $38$ of the $\binom{8}{4}=70$ four-element subsets are maximal.
\end{itemize}

\subsection{Complexity and observed costs}

Table~\ref{tab:cost} summarises the routines.  All orientation predicates are exact
rational arithmetic, so the bit-complexity depends on the denominator growth of the
inputs; for the integer inputs of the order-type database this is negligible.

\begin{table}[ht]
\centering
\caption{Cost of the main routines.  $\Sigma$ is the number of nodes visited and
$d$ the maximum degree of the forbidden-quadruple hypergraph, which is
$O(n^3)$.}
\label{tab:cost}
\small
\begin{tabular}{@{}lll@{}}
\toprule
routine & time & worst case \\
\midrule
$\operatorname{orient}$ & $O(1)$ & exact rational arithmetic \\
$\conv$, hull of $k$ points & $O(k\log k)$ & monotone chain \\
\texttt{maximal\_convex\_\allowbreak subsets}
  & $O(2^n n\log n)$ & the $2^n$ subset scan \\
\texttt{maximal\_convex\_\allowbreak subsets\_fast}
  & $O(n^3+\Sigma n d)$ & output-sensitive \\
  & & $O(n|\M(P)|)$ nodes \\
\texttt{enumerate\_order\_\allowbreak types\_up\_to}
  & last level dominates & complete $n\le 7$; \\
  & & $3313/3315$ at $n=8$ \\
\bottomrule
\end{tabular}
\end{table}

Here $\Sigma$ is the number of nodes visited and $d$ the maximum degree of the
forbidden-quadruple hypergraph, which is $O(n^3)$.  The important point is that the
second routine is exponential in the \emph{output} $|\M(P)|$ rather than in $2^n$.

On a single core, scanning all order types takes $4$ s for $n=7$, $132$ s for
$n=8$, and $18\,483$ s ($\approx 5.1$ h) for $n=9$.  Extrapolating, $n=10$
would need $\approx 19$ days, which is why Open Problem 8.2 stands.  The
output-sensitive routine exhausts a $4\times 10^8$-node budget on a $16$-point
twin-pair configuration and raises; for $n\le 14$ it completes.  This is the
intended behaviour --- it reports failure rather than returning a truncated count.

\section{Constructions and falsification}
\label{sec:twin}

\subsection{Twin pairs, and a caution about lower bounds}

Take $k$ directions and, in each, two nearly coincident points at radii $1$ and
$1-\varepsilon$; this gives $2k$ points.  One is tempted to expect that, as
$\varepsilon\to 0$, the configuration degenerates to a ``twin pair'' pattern in
which a maximal convex subset picks one point from each pair, giving $2^k$.

\begin{table}[ht]
\centering
\caption{$|\M(P)|$ for twin-pair configurations with $n=2k$ points, computed
exactly.  Compare the last two columns with the exhaustive values
$4,11,20,38,62$ for $f(4),\dots,f(8)$: the construction is \emph{weaker} than
exhaustive search for small $n$, and its values are not monotone in $k$ for
fixed $\varepsilon$.}
\label{tab:twin}
\begin{tabular}{c c r r r r r r r}
\toprule
$k$ & $n$ & $\varepsilon=\frac12$ & $\frac14$ & $\frac18$ & $\frac1{16}$ &
$\frac1{32}$ & $\frac1{64}$ & $\frac1{128}$ \\
\midrule
3 & 6  & 4  & 4  & 4  & 4  & 4  & 4  & 4 \\
4 & 8  & 7  & 11 & 14 & 11 & 11 & 11 & 11 \\
5 & 10 & 17 & 21 & 24 & 24 & 26 & 26 & 26 \\
6 & 12 & 26 & 36 & 55 & 49 & 51 & 57 & 57 \\
7 & 14 & 34 & 53 & 71 & 93 & 113 & 121 & 120 \\
\bottomrule
\end{tabular}
\end{table}

Table~\ref{tab:twin} needs an honest reading.  Each entry is
$|\M(P)|$ for one \emph{specific} configuration, hence a lower bound for $f(2k)$
--- but a weak one.  At $n=8$ the best separation gives $14$, while exhaustive
search over all $3315$ order types gives $f(8)=38$.  So for small $n$ the
construction is far from optimal, and we explicitly withdraw the expectation
$2^{n/2}$ that motivated it: the twin-pair limit is evidently not the mechanism
behind the true optimum, which is attained by configurations of a quite different
shape (all maximal subsets of size exactly $n/2$, with equal point-degrees).

What the table does establish is that $|\M(P)|$ is \emph{extremely} sensitive to
the exact order type: on fixed $14$ points, varying only a rational separation
parameter takes $|\M(P)|$ from $34$ to $121$, a factor of $3.5$.  Any
asymptotic argument for $f$ must therefore control the order type, not just the
point count.

\subsection{Falsification attempts}

We tried hard to break every theorem and record the failures honestly.

\paragraph{Duality (Theorem~\ref{thm:duality}).} Verified on all $3472$ order types
with $n\le 8$ and on $400$ random configurations for each $n=9,\dots,16$.  No
counterexample found.

\paragraph{Kill lemma (Lemma~\ref{lem:kill}).} Verified on $6126$ pairs $(S,p)$ with
$n\le 7$ and $p\notin\conv(S)$; zero violations.

\paragraph{A rejected conjecture.} We conjectured that every maximal convex subset
$S$ satisfies $|S|\ge\min(h,n-h)$ where $h$ is the hull size, i.e.\ that maximisers
live near $h\approx n/2$.  This is \emph{false}: Table~\ref{tab:layer} shows the
maximum at $h=3$ or $4$, and we found $1292$ violations for $n\le 8$.  We discard
it.

\paragraph{A rejected injection.} The map $S\mapsto S\setminus\{\text{leftmost point of } S\}$
is \emph{not} injective (it fails already at $n=4$), so a bound
$|\M(P)|\le 2^{n-2}$ is not obtainable this way.  The map $\psi$ of
Theorem~\ref{thm:duality} \emph{is} injective, but its image is far from all
subsets, so it does not by itself sharpen Sperner.

\paragraph{A rejected expectation.} We expected the twin-pair construction to
recover $2^{n/2}$ in the limit $\varepsilon\to0$.  Table~\ref{tab:twin} refutes
this: the construction's best value at $n=8$ is $14$, far below the exhaustive
$f(8)=38$.  The mechanism behind the true optimum is different from the one we
assumed, and we say so rather than quietly dropping the experiment.

\section{Discussion and open problems}
\label{sec:open}

\paragraph{What we know.} $f$ is a function of $n$ only; $|\M(P)|$ is order-type
invariant; $|\M(P)|\le\binom{n}{\lfloor n/2\rfloor}$ by Sperner; the values
$f(3),\dots,f(9)$ are as in Theorem~\ref{thm:f}.

\begin{conjecture}[growth of $f$]
\label{conj:growth}
$f(n)=2^{\Theta(n)}$.
\end{conjecture}

\begingroup
\sloppy
\noindent
\emph{Evidence, and its limits.}  The exact values $1,4,7,11,20,38,62$ have
successive ratios $4.0,1.75,1.57,1.82,1.90,1.63$, and $f(7)=20$, $f(8)=38$,
$f(9)=62$ are well fitted by $c$ times $1.87^n$.  The data are consistent with
exponential growth, but are far too short to distinguish $c$ times $2^{n/2}$
from $c$ times $1.87^n$.  Against this, our explicit constructions
(Section~\ref{sec:twin}) are weaker than the exhaustive values already at $n=8$,
so they give no useful asymptotic lower bound, and we have \emph{no} construction
we can prove beats $2^{n/2}$.  On the upper side we have only Sperner.  We
therefore regard Conjecture~\ref{conj:growth} as suggestive rather than well
supported, and we are explicit that we prove neither direction.
\par
\endgroup

\begin{openproblem}
Determine $f(10)$.  The database has $14\,309\,547$ order types of $10$ points
($\approx 572$ MB), so this is a straightforward but expensive computation, which we
did not run.
\end{openproblem}

\begin{openproblem}
Characterise the maximisers.  At $n=8$ the maximisers have all maximal subsets of
size $n/2$ and all point-degrees equal; we do not know whether this balanced
structure is forced for large $n$ or is merely typical of small instances.
\end{openproblem}

\begin{openproblem}
Determine the largest possible number of maximal convex subsets of size exactly
$3$.  Such a subset is a triangle $T$ with $P\setminus T\subseteq\int(\conv(T))$, a
strong condition that our data suggests makes them rare ($6$ at $n=5$, $2$ at
$n=6$, $1$ at $n=7$).
\end{openproblem}

\paragraph{Limitations.} (i) Our independent order-type enumerator is incomplete
at $n=8$ ($3313/3315$ types); the exact values rely on the external database,
which we validated but did not produce.  (ii) The upper bound is Sperner's and is
far from sharp.  (iii) The constructions in Table~\ref{tab:twin} are unoptimised;
we used simulated annealing on coordinates only as a screening tool and do not
claim its outputs are optimal.  (iv) No asymptotics are proved.

\paragraph{Honest score.} We regard the exact values for $n\le 9$, the duality
theorem and the reproducible computational pipeline as the strongest parts of this
work.  The asymptotic question is open, the upper bound is weak, and the paper is
a modest but honest computational-structural contribution rather than a
breakthrough.

\begin{thebibliography}{9}

\bibitem{baranyvaltr04} I. B\'ar\'any and P. Valtr. Planar point sets with a small number
of empty convex polygons. \emph{Studia Sci.\ Math.\ Hungar.\ }41(2):243--266, 2004.

\bibitem{bae2022} S.~W. Bae. Faster counting empty convex polygons in a planar point
set. \emph{Inf.\ Process.\ Lett.\ }175:106221, 2022.

\bibitem{erdos838} Erd{\\H{o}}s Problem 838 (T.~Bloom, ed.). Minimum number of convex subsets
determined by $n$ points. \url{https://www.erdosproblems.com/838}.

\bibitem{hmp2020} A.~F. Holmsen, H.~Nassajian Mojarrad, J.~Pach, and G.~Tardos. Two
extensions of the Erd{\\H{o}}s--Szekeres problem. \emph{J.\ Eur.\ Math.\ Soc.\ }22:3741--3763, 2020.

\bibitem{bukhmnk11} A.~Bukh, J.~Matou\v{s}ek, and G.~Nivasch. Lower bounds for weak
$\varepsilon$-nets and stair-convexity. \emph{Israel J.\ Math.\ }182(1):199--228, 2011.
(Conference version: SoCG 2009, 1--10.)

\bibitem{huemer2022} C.~Huemer, D.~Oliveros, P.~P\'erez-Lantero, F.~Torra, and
B.~Vogtenhuber. On weighted sums of numbers of convex polygons in point sets.
\emph{Discrete Comput.\ Geom.\ }68, 2022.

\bibitem{mitchell95} J.~S.~B. Mitchell, G.~Rote, G.~Sundaram, and G.~Woeginger.
Counting convex polygons in planar point sets. \emph{Inf.\ Process.\ Lett.\ }
56(1):45--49, 1995.

\bibitem{pinchasi2006} R.~Pinchasi, M.~Radi\'ci\'c, and M.~Sharir. On empty convex
polygons in a planar point set. \emph{J.\ Combin.\ Theory Ser.\ A }113(3):385--419, 2006.

\bibitem{rote91} G.~Rote, G.~Woeginger, B.~Zhu, and Z.~Wang. Counting $k$-subsets and
convex $k$-gons in the plane. \emph{Inf.\ Process.\ Lett.\ }38(3):149--151, 1991.

\bibitem{rote92} G.~Rote and G.~Woeginger. Counting convex $k$-gons in planar point
sets. \emph{Inf.\ Process.\ Lett.\ }41(3):191--194, 1992.

\bibitem{savova2026} N.~Savova. Point sets with few convex subsets: iterated blow-ups
and exact small values. Preprint, 2026.
\url{https://github.com/NikolSavova/erdos-838-convex-subsets}.
(Not peer reviewed; cited as adjacent, competing work.)

\bibitem{suk2017} A.~Suk. On the Erd{\\H{o}}s--Szekeres convex polygon problem.
\emph{J.\ Amer.\ Math.\ Soc.\ }30(4):1047--1053, 2017.

\bibitem{aak2002} O.~Aichholzer, F.~Aurenhammer, and H.~Krasser. Enumerating order types
for small point sets with applications. \emph{Order }19(3):265--281, 2002.

\bibitem{aak2001} O.~Aichholzer and H.~Krasser. The point set order type data base: a
collection of applications and results. \emph{CCCG 2001}, 17--20.

\bibitem{aak2007} O.~Aichholzer and H.~Krasser. Abstract order type extension and new
results on the rectilinear crossing number. \emph{Comput.\ Geom.\ }36(1):2--15, 2007.

\bibitem{aaDB} Aichholzer order-type database.
\url{http://www.ist.tugraz.at/staff/aichholzer/research/rp/triangulations/}
\allowbreak
\url{ordertypes/} (accessed October 2026).

\end{thebibliography}

\end{document}